\section{La conception d'un système politique idéal}

\subsection*{Introduction}

\subsection{Les fondements naturels et républicain}

\subsubsection{L'utopie d'un gouvernement du peuple}

Ambroise de Milan ne se distingue pas parmi les Pères de l'Église comme un auteur majeur de théorie politique comme peut l'être Augustin, et pour cause aucune de ses œuvres ne s'attelle à réfléchir avec précision à une nouvelle conception du pouvoir ou de la diplomatie romaine. Pourtant, il est bien possible en s'attaquant à la lecture de l'évêque de trouver une théorisation poussée d'un idéal politique original. C'est principalement le cas du livre V de son ouvrage \latin{Exameron} qui comporte des bribes majeures de la pensée politique ambrosienne. Cette œuvre théologique est un recueil d'homélies prononcées durant la Semaine Sainte, probablement à la fin des années 380, qui propose un commentaire détaillé des six jours de la Création. Ambroise s'inspire pour cette prédication des sermons de Basile de Césarée et d'Origène sur ce même thème\autocite[100-104]{fusco2025}, en y ajoutant donc entre autres une critique politique de la société humaine, ce qui nous permet de saisir sa pensée sur le pouvoir romain.

\bigskip

Cette pensée politique justement se construit derrière l'une des idées centrales chez Ambroise : une répartition absolue des tâches, des droits et des devoirs, et surtout du pouvoir pour aboutir à un « gouvernement parfait ». Pour concevoir et argumenter sa propre réflexion, Ambroise ne se tourne pas vers les traités juridiques ou les constitutions romaines, mais vers l'observation de la nature afin de tirer profit d'une analyse fine de ce qu'il considère être l'œuvre directe de Dieu. En décrivant le comportement des animaux, Ambroise dresse en creux le portrait des institutions humaines idéales, car pour lui la nature originelle est le reflet de l'ordre voulu par Dieu. Son observation débute avec les oiseaux, et notamment l'organisation sociale des grues où il observe un fonctionnement dicté par la nature qu'il oppose aux sociétés humaines :

\bigskip

\begin{quote}
    «~Commençons donc par les êtres qui nous ont offert, pour notre usage, un modèle à imiter. Chez eux, en effet, on trouve une sorte de vie politique et une organisation militaire naturelles, alors que chez nous, elles sont contraintes et serviles\footnote{« \latin{Ab his igitur ordiamur quae nostro se usui imitationem dederunt. in illis enim politia quaedam et militia naturalis, in nobis coacta atque seruilis}.». \cite[Exameron, V, 15, 50]{ambroise_csel_32_1}.}~»
\end{quote}

\bigskip

Dans son immense ouvrage sur la pensée politique chrétienne et byzantine, l'historien Francis Dvornik s'arrête longuement sur la philosophie ambrosienne et sur son impact dans la théologie politique, ce qui en fait un passage majeur dans le cadre de ce mémoire. Il observe donc tout particulièrement son avis sur la monarchie à travers les sociétés naturelles, et analyse alors l'importance de cette «~\latin{politia quaedam et militia naturalis}~» dans la pensée d'Ambroise\autocite[674]{dvornik_political}. Cette formule est ce qui reflète le mieux l'idée que la nature n'est pas seulement un décor mais bien un exemple d'une organisation sociale idéale avec une vision ordonnée du monde où chacun a une place et des responsabilités. Il est intéressant de noter que dans la complexité de son analyse, Dvornik associe cet ordre de la nature aux théories politiques des monarchies hellénistiques, affirmant que la hiérarchie et la discipline sont issues d'un ordre divin. C'est une des critiques qu'il est possible de faire sur l'œuvre de Dvornik lorsqu'il relie de nombreuses réflexions chrétiennes à des théories valorisant la royauté : «~La société idéale d'Ambroise reflète à nouveau l'État hellénistique : un système d'ordre et de répartition du travail [...]\autocite[676]{dvornik_political}~». En effet, bien que comme nous le verrons dans la suite de ce chapitre, Ambroise ne s'oppose pas frontalement à l'ordre monarchique et est même bien loin de se révolter contre les souverains, son passage sur l'ordre naturel et notamment l'exemple des grues valorise un modèle collégial et républicain, comme l'indique la citation qui a ouvert ce mémoire :

\bigskip

\begin{quote}
    «~Quoi de plus beau que de voir que le travail et l’honneur sont communs à tous, et que le pouvoir n’est pas un privilège de quelques-uns mais qu’il est réparti de façon volontaire entre tous ?\footnote{« \latin{Quid hoc pulchrius, et laborem omnibus et honorem esse communem nec paucis adrogari potentiam, sed quadam in omnes voluntaria sorte transcribi ?}.». \cite[\nopp V, 15, 51]{ambroise_csel_32_1}.}~»
\end{quote}

\bigskip

Chez les grues, Ambroise loue un service social et militaire où le travail est constant et n'est jamais remis en question\footnote{Voir le 1.1.2 pour le développement de cette idée.} et plus encore, où le rôle de dominant n'est pas figé dans les mains d'un seul : dans les formations de vol, l'oiseau placé en tête se fatigue pour guider la communauté avant d'être volontairement remplacé par un autre. Ce roulement assure confiance et efficacité dans la vie en société. Dans ce passage, Ambroise projette son admiration pour cette société solidaire en formulant une sorte de plaidoyer pour le partage du pouvoir. C'est ici le cœur de son utopie politique, une forme de paroxysme de la \latin{Res publica}, cette «~chose publique~» qui donne au peuple la gestion de la communauté et de la vie de la cité dans son ensemble, sans prise de pouvoir exclusif. L'évêque affirme donc une volonté républicaine forte et sans ambiguïté, où le pouvoir ne doit pas être un privilège, mais doit se soumettre à une rotation volontaire. Cet idéal ambrosien permet à la fois une répartition du pouvoir, mais également une lutte contre la crainte de la servitude, puisqu'elle est partagée au même titre que le pouvoir. Les contraintes sont toujours temporaires et compensées par l'accès de chacun, à un moment donné, au pouvoir et à l'honneur.

\bigskip

Et l'un des passages les plus intéressants de cette réflexion est lorsque l'évêque de Milan transpose cette société du monde animal à l'histoire de l'être humain, affirmant notamment que l'Homme a perdu cet idéal, mais qu'il a été une réalité concrète des premières communautés :

\bigskip

\begin{quote}
    «~C’est ainsi qu’à l’origine, les hommes, suivant l’exemple des oiseaux, avaient commencé à exercer une vie politique reçue de la nature : le travail était commun, la dignité commune ; chacun apprenait à partager les soucis à tour de rôle, à diviser les obéissances et les commandements, de sorte que nul ne fût exclu de l'honneur, et nul exempt du travail.\footnote{« \latin{antiquae hoc rei publicae munus et instar liberae ciuitatis est. sic a principio acceptam a natura exemplo auium politiam homines exercere coeperant, ut communis esset labor, communis dignitas, per uices singuli partiri curas discerent, obsequia imperiaque diuiderent, nemo esset honoris exsors, nullus inmunis laboris}.~». \cite[\nopp V, 15, 52]{ambroise_csel_32_1}.}~»
\end{quote}

\bigskip

À travers ce texte, Ambroise dresse d'abord le tableau d'une humanité ayant chuté, rappelant par là le Péché originel, et qui a perdu sa nature première d'égalité et de rejet de la domination. L'idée de l'évêque est de démontrer que son idéal du gouvernement et de la vie en société n'est pas un modèle irréalisable et utopique de la \latin{Res publica}, mais est plutôt la nature première de l'être humain, celle de l'ordre naturel institué par Dieu et qui aujourd'hui ne se voit plus que dans certaines sociétés animales comme les grues. L'aboutissement à des systèmes monarchiques et impériaux est donc selon Ambroise le résultat d'un échec de la sauvegarde de la nature humaine. Une grande partie de son idéal politique et de sa critique des sociétés vient de cette idée que les Hommes ont perdu leur capacité naturelle à suivre l'ordre établi par Dieu, corrompus par l'appât du gain, la soif de pouvoir et la sauvegarde de la propriété privée. C'est donc devant cette incapacité de la communauté de rendre des services par essence et par charité que l'humanité a dû se doter d'institutions strictes et exigeantes. La perte de cet état de nature semble donc justifier, par défaut, l'existence d'un pouvoir impérial contraignant, bien loin de l'utopie ambrosienne de la République.

\subsubsection{Une société d’entraide pour faciliter le bon gouvernement}

L'utopie de la \latin{Res publica} et d'un pouvoir véritablement partagé, tel qu'Ambroise le théorise à travers les exemples de la nature, ne constitue en soi qu'une petite partie de sa pensée politique et de sa conception d'un système politique idéal. L'évêque de Milan est bien conscient que cette demande de répartition des responsabilités et des privilèges ne peut pas sortir à partir de rien alors que le monde romain est ancré depuis longtemps maintenant sous la domination d'un ou de plusieurs empereurs. Dans ses quelques textes qui parlent ouvertement de sa vision de la société ou du pouvoir, Ambroise rappelle régulièrement qu'un bon gouvernement, du moins ce qu'il considère comme tel, repose tout autant sur l'attitude des gouvernants que des gouvernés. Autrement dit, dans l'objectif assez limpide d'aboutir à un « pouvoir réparti de façon volontaire entre tous » pour éviter toute approche tyrannique, il est nécessaire d'avoir une implication active de la bonne volonté du peuple. Toute constitution conforme à la nature ne peut exister que par le libre accord et l'entraide constante entre les membres de la société, quels qu'ils soient. Le questionnement autour de l'entité détenant le pouvoir est donc nécessaire, nous le verrons notamment dans la suite de ce chapitre sur l'acceptation de la monarchie, mais les actions et implications de la population le sont tout autant.

\bigskip

Comme le souligne l'historien Giuseppe Zecchini, Ambroise n'utilise pas ses écrits pour remettre en cause le principe monarchique ou impérial de son temps. Il en fait même pleinement partie, que ce soit en tant que gouverneur de la province de Ligurie-Émilie jusqu'en 374, puis en tant qu'évêque de la capitale impériale et ambassadeur officiel de l'empereur Valentinien II auprès de l'usurpateur Maxime entre 383 et 387. Ambroise ne se présente donc pas comme révolutionnaire, mais plutôt comme un penseur capable d'influencer la société pour la rendre plus juste. Dans ce cadre, il a conscience de la nécessité d'entretenir un lien important entre la population et les détenteurs du pouvoir. C'est ce que Zecchini appelle « concilier la \latin{parcitas imperiale} (la retenue de l'empereur) avec la \latin{libertas} des citoyens\footnote{« \latin{Ora, Ambrogio non discute il principio monarchico, ma ha ben presente l'esigenza di conciliare la parcitas imperiale con la libertas dei cittadini}.~». \cite[164]{zecchini_pensiero_politico}.}. Cette retenue s'exprime principalement par l'exigence de la foi et de la morale chrétienne que l'évêque de Milan n'hésite pas à rappeler dans chacun de ses passages de théorie politique en parlant de «~zèle~» y compris chez les animaux. En revanche la liberté civique dont parle Zecchini et qui est chère à Ambroise n'est pas développée comme un droit passif : elle n'est accessible que par la participation de tous aux devoirs de la communauté. L'évêque reprend ainsi l'exemple de la société des grues de son \latin{Exameron} pour illustrer ce libre accord qui concorde avec une lutte contre la paresse. Il s'adresse à la société romaine en décrivant la grue qui prend son tour de garde :

\bigskip

\begin{quote}
    «~Cette dernière accepte son sort de bon gré et, contrairement à notre habitude, elle ne renonce pas au sommeil à contrecœur ou avec paresse ; au contraire, elle s'arrache vivement à sa couche, accomplit son tour et rend, avec un soin et un zèle égaux, le service qu'elle a reçu.\footnote{« \latin{At illa uolens suscipit sortem nec usu nostro inuita et pigrior somno renuntiat, sed inpigre suis excutitur stratis, uicem exsequitur et quam accepit gratiam pari cura atque officio repraesentat}.~». \cite[\nopp V, 15, 50]{ambroise_csel_32_1}.}~»
\end{quote}

\bigskip

Ambroise dresse ici une critique de la paresse inhérente aux sociétés humaines qui, dans sa pensée, est un facteur entretenant l'accaparement du pouvoir et l'injustice générale. Le bon gouvernement dépend d'une confiance mutuelle qui ne peut être instaurée que si chacun accomplit ses devoirs, même les plus fastidieux, avec la même volonté. Pour aboutir à une répartition du travail, l'évêque en appelle dans un premier temps à la répartition du travail de tous sans quoi les privilèges, les envies privées et la jalousie ne peuvent être combattus. La répartition du pouvoir dont il est question devient donc viable par l'accompagnement d'une acceptation volontaire des contraintes, qui là encore est quelque chose de naturel puisqu'on le retrouve dans le monde animal. Il y a donc toujours une critique de l'humain dans sa globalité face à l'ordre naturel établi par les Six jours de la Création dont Ambroise fait le commentaire. À travers cette vision, Ambroise semble même sous-entendre qu'un gouvernement est presque forcé de devenir autoritaire pour imposer par la contrainte les devoirs qui devraient être faits naturellement. Il évoque à ce sujet le fait que l'homme, par son refus de l'accomplissement des devoirs, oblige la venue d'un gouvernement coercitif afin de réaliser les tâches uniquement sous la menace : il parle de «~peine pour la paresse\footnote{« \latin{proponitur poena desidiae}.~». \cite[\nopp V, 15, 52]{ambroise_csel_32_1}.}~» comme le signe d'un échec du travail naturel.

\bigskip

Pour approfondir cette vertu politique du travail et du dévouement à la communauté, Ambroise mobilise un peu plus loin dans son \latin{Exameron} une deuxième figure du monde animal : l'abeille. Cet exemple tout aussi important dans sa pensée politique lui sert principalement pour exprimer concrètement sa vision de la monarchie et le réalisme qu'il porte à son idéal de la société\footnote{Cette idée est développée à partir de la sous-partie 2.1.1}, mais il est ici utilisé pour parler du travail en société et de l'importance du bien commun. Chaque individu connaît sa tâche et s'y attelle sans contestation :

\bigskip

\begin{quote}
    «~Va vers l’abeille et vois comme elle est active, comme elle est admirable dans son activité, dont les résultats du travail sont utiles aux rois et aux gens du commun ; elle est désirée par tous et illustre.\footnote{« \latin{Uade ad apem et uide quomodo operaria est, operationem quoquequam uenerabilem mercatur, cuius laborem reges et mediocres ad salutem sumunt ; adpetibilis enim est omnibus et clara}.~». \cite[\nopp V, 21, 70]{ambroise_csel_32_1}.}~»
\end{quote}

\bigskip

Ce passage est évidemment relié dans la réflexion d'Ambroise au passage biblique suivant : «~Va vers la fourmi, paresseux, considère ses voies, et tu deviendras sage : elle qui n’a point de chef, ni de surveillance, ni de maître, elle se procure pendant l’été son pain , elle amasse pendant la moisson sa nourriture\footnote{Proverbes, 6, 6-8.}~». Bien que la comparaison ne soit pas identique, l'idée qu'elle porte reste la même : démontrer que le travail au service de la communauté est le seul moyen de garantir une vie sociale, et donc d'aboutir à une véritable \latin{Res publica} ne nécessitant plus de souverain unique. Il est ici question de valorisation de toute forme de travail comme quelque chose profitant à tous, aussi bien «~aux souverains qu'aux gens du commun~», et donc comme un élément primordial d'une société juste aboutissant à un bon gouvernement. Si on reprend la comparaison à la citation du livre des Proverbes, cette lutte contre la paresse humaine est ce qui permet de se retrouver sans «~chef, ni surveillance, ni maître~», ce qui explique l'importance que met Ambroise dans la réalisation des devoirs. L'activité et l'entraide gomment ainsi les hiérarchies face aux nécessités du bien public, seul moyen d'aboutir à un système républicain.

\bigskip

La réflexion théorique ambrosienne sur les meilleurs moyens d'aboutir à un système politique idéal ne se limite pas aux sermons sur la Création. Elle trouve également sa place dans son \latin{De Officiis}, un traité à l'attention du clergé, grandement calqué sur le \latin{De Officiis} de Cicéron, dont la rédaction prend fin vers 388-389. Ambroise y développe de nombreuses idées d'éthique chrétienne et de théologie morale que son clergé doit appliquer, faisant de l'entraide un point incontournable de la société et du bon gouvernement :

\bigskip

\begin{quote}
 «~Ainsi donc selon la volonté de Dieu ou le lien de la nature, nous devons nous secourir mutuellement, rivaliser dans l’accomplissement des devoirs, mettre pour ainsi dire en commun tous les intérêts et, pour user du mot de l’Écriture, nous porter assistance l’un à l’autre, ou bien par le zèle, ou bien par l’accomplissement du devoir, ou bien par l’argent, ou bien par les œuvres, ou bien de n’importe quelle manière, afin d’accroître l’agrément du lien social entre nous\footnote{« \latin{texte en latin}.~». \cite[\nopp I, 135]{ambroise_devoirs_1}.}.~»
\end{quote}

\bigskip

Ambroise réaffirme dans ce passage l'idée que le travail en commun et la vie en communauté dans sa globalité viennent de l'ordre naturel des choses et qu'il est donc nécessaire d'entretenir ce mode de fonctionnement pour le bien de tous. Mais l'évêque de Milan ajoute ici une énumération fondamentale de sa théorie politique sur l'entraide. Il ne fait plus seulement une comparaison avec le monde animal, mais évoque et décrit les modalités concrètes du partage. On y retrouve en premier lieu ce zèle que nous avons évoqué : il n'oublie pas sa position d'évêque et veut avant tout montrer que la foi chrétienne est le meilleur moyen d'aboutir à son idéal politique. Puis il est question d'une dimension civique avec l'accomplissement du devoir, ce qui montre une certaine continuité dans sa pensée entre la valorisation du travail des abeilles et ce passage : la république est un modèle demandant une véritable implication de ses citoyens. Enfin il dresse des faits plus concrets en parlant d'aide «~par l'argent ou bien par les œuvres~». On comprend bien qu'à travers ce passage du \latin{De Officiis}, Ambroise ne conçoit pas une théorie politique qui demande un changement de régime, mais bien une réflexion sur l'ensemble de la société afin de trouver un équilibre entre la vie en communauté et le bon gouvernement.

\bigskip

La façon qu'a l'évêque de Milan de décrire l'état naturel de l'Homme et de l'animal dans son \latin{Exameron} ne laisse pas de doute quant à son appréciation d'un pouvoir partagé. Mais à travers ses écrits, on comprend surtout qu'il place cet idéal dans un coin presque inatteignable de sa pensée, et qu'il utilise cette finalité comme un argument pour pousser la société romaine au changement, à commencer par la lutte contre la paresse de la population. La réalisation d'un «~lien social~» n'est pas seulement une demande chrétienne, c'est un argument crucial pour la \latin{Res publica}. En incluant dans l'entraide communautaire aussi bien les détenteurs du pouvoir, peu importe sa forme, que le peuple, il pose les bases d'un système politique dont la réussite se mesure par la solidarité de ses citoyens.

\subsubsection{La critique de la soif de pouvoir}

Le rejet de la paresse et la demande explicite d’un effort communautaire de travail, tels que nous venons de les analyser, constituent le socle d'une société plus juste, du moins dans la vision qu'en a Ambroise. Cependant, cette demande d'un effort collectif, qui peut véritablement être vu comme requête concrète puisqu'il s'agit d'un sermon, pour modeler le bon gouvernement depuis un mouvement de la communauté et non pas seulement par un changement des détenteurs du pouvoir, doit s'accompagner dans le discours de l'évêque d'une critique stricte du pouvoir en lui-même, et donc de la recherche de celui-ci. En effet, la lecture des textes d'Ambroise, et de nouveau principalement son \latin{Exameron}, met en lumière un aspect fortement décrié et qui s'inscrit au coeur de l'échec de l'état naturel voulu par la Création : le désir innarétable de pouvoir qui caractérise l'être humain. D'uen façon générale, la pensée politique ambrosienne se construit sur la critique d'un accaparement des privilèges et sur le rejet d'une revendication du pouvoir et de l'autorité par la force. Dans ce cadre, il semble normal qu'il condamne la conservation d'un commandement dans le temps long, le pointant du doigt comme une faute morale. C'ets dans le paragraphe 52 du livre V de l'\latin{Exameron} qu'il développe le plus clairement ces idées et critiques, venant définir avec précision son idéal politique fondé sur des modèles naturels et républicains que nous avons précédemment évoqué.

\bigskip

Ambroise introduit son passage critique sur la confiscation du pouvoir en résumant l'idéal naturel qu'il a observé chez les grues par une formule impactante :

\bigskip

\begin{quote}
    «~C'est là la fonction d'une ancienne république et l'image d'une cité libre.\footnote{« \latin{Antiquae hoc rei publicae munus et instar liberae ciuitatis est}.~». \cite[\nopp V, 15, 52]{ambroise_csel_32_1}.}~»
\end{quote}

\bigskip

Cette phrase est un condensé de la pensée ambrosienne sur la comparaison entre son idéal politique et l'état des sociétés naturelles. Il revient à la fois sur le travail pour le bien commun et sur le partage des responsabilités pour décrire ce qui correspond pratiquement à son aboutissement théorique. Il est intéressant d'analyser cette citation car l'évêque de Milan y utilise des formules typique de la pensée politique et juridique romaine, pour jutsemment décrire un modèle qui difère justemment grandement de Rome. On y retrouve une notion que nous avons déja utilisé dans ce mémoire : la \latin{res publica} ainsi que celle de \latin{libera civitas}. Il analyse donc d'autres sociétés en les comparant à des concepts romains, opérant ainsi une sorte de rejet implicite du fonctionnement politique contemporain. Le fait de faire raisonner l'idée de liberté est ici d'autant plus crucial car elle permet de faire voir l'accaparement du pouvoir comme un modèle de domination et de servitude.

\bigskip

Ambroise continue alors son sermon par l'un des passages les plus importants de cette réflexion, décrivant la prise et la conservation du pouvoir comme un élément non seulement contre nature, mais surtout rendant impossible l'utopie d'un gouvernement agissant pour le bien commun :

\bigskip

\begin{quote}
    «~Mais depuis que la passion de dominer a commencé à revendiquer les pouvoirs obtenus et à ne plus vouloir déposer ceux dont elle s'est saisie ; depuis que la vie militaire a cessé d'être un droit commun pour devenir une servitude ; depuis que l'accession au pouvoir n'est plus un ordre établi mais une ambition de s'en emparer, l'accomplissement même du travail a commencé à être supporté plus durement, et ce qui n'est pas volontaire laisse vite place à la négligence.\footnote{« \latin{Sed postquam dominandi libido uindicare coepit indeptas et susceptas nolle deponere potestates, posteaquam militiae non ius commune coepit esse, sed seruitus, posteaquam non ordo factus est suscipiendae potentiae, sed studium uindicandae, coepit etiam ipsa laboris functio durius sustineri, et quae non est uoluntaria cito locum relinquit incuriae}.». \cite[\nopp V, 15, 52]{ambroise_csel_32_1}.}~»
\end{quote}

\bigskip

Ambroise dépose dans ce passage une critique ferme des ambitions privées de l'Homme, qu'il désigne comme l'élément primaire de l'incapacité des sociétés républicaines à pleinement se développer entre les mains du peuple. On note d'ailleurs que le passage d'un ordre naturel à une ambition personnelle est relié à la recherche du travail et à la réalisation des devoirs, ce qui montre que la réalisation de son idéal ne se fait que par un movement généralisé de l'ensemble des aspects impactants un gouvenrement. Dans cet extrait, le vocabulaire utilisé se révèle tout aussi important avec le terme «~\latin{dominandi libido}~» qui fait de la prise de pouvoir et du sentiment d'autorité un désir, une passion, propre à l'être humain contemporain, et donc un désir qu'Ambroise veut combattre puisque qu'il rentre en contradiction avec la forme première de la Création. La confiscation du pouvoir par un seul ou par un groupe apparaît dans la pensée ambrosienne comme une sorte de péché originel de la vie politique, établissant à jamais des régimes monarchiques et tyranniques. Le pouvoir qui s'éternise dans les mains d'un seul corrompt aussi bien l'homme qui le détient que la communauté sous ses ordres. L'idée de servitude revient une nouvelle fois comme la conclusion inévitable pour une population vivant dans l'ombre de l'ambition personnelle et de la recherche du pouvoir.

\bigskip

Enfin, un peu plus loin dans ce long passage sur la critique humaine du désir du pouvoir, Ambroise cherche à confirmer que sa critique n'est pas qu'un simple aspect théorique mais bien une réalité contemporaine concrète qui influe la psychologie humaine :

\bigskip

\begin{quote}
    «~C’est pourquoi un travail acharné et constant détourne l'affection, et une puissance continue et durable engendre l'insolence. Quel homme trouverais-tu qui abandonnerait de lui-même le pouvoir, renoncerait aux insignes de son commandement et accepterait de passer volontairement du premier rang au dernier ? Quant à nous, non seulement nous luttons souvent pour la première place, mais même pour une place intermédiaire ; nous revendiquons les premières places dans les banquets et, si cet honneur nous est accordé une fois, nous voulons qu'il soit perpétuel.\footnote{« \latin{Ergo et iugis labor auertit affectum et continua ac diuturna potentia gignit insolentiam. quem inuenias hominum, qui sponte deponat imperium et ducatus sui cedat insigne fiatque uolens numero postremus ex primo! nos autem non solum de primo, sed etiam de medio saepe contendimus et primos discubitus in conuiuio uindicamus aec, si semel delatum fuerit, uolumus esse perpetuum}.~». \cite[\nopp V, 15, 52]{ambroise_csel_32_1}.}~»
\end{quote}

\bigskip

L'évêque de Milan pose dans ce paragraphe une sorte de constat de la réalité de la société romaine. Son rejet théorique de la monarchie, puisque nous verrons qu'il se montre en réalité très réaliste face au pouvoir impérial et qu'il vise un changement dans la gestion du pouvoir plutôt qu'un renversement, se résume particulièrement bien ici : «~\latin{continua ac diuturna potentia gignit insolentiam}~», «~une puissance continue et durable engendre l'insolence~». Ambroise appuie sur le fait que cet échec est lié à l'ensemble des citoyens qui cherchent à défendre leur pouvoir, aussi futile soit-il, plutôt qu'oeuvrer pour le bien de la communauté et aboutir à une véritable \latin{Res publica}. L'acharnement envers la lutte et la recherche du pouvoir ne se fait pas que dans le cadre du pouvoir ou de l'autorité la plus grande, mais bien dans tout les aspects de la société où chacun logne pour la récupération d'un pouvoir toujours plus grand. L'év6eque de Milan démontre une fois encore que l'aboutissement à des régimes coerciftifs et tyranniques n'est que la conséquence d'un échec général de la société humaine. L'état naturel de l'homme semble donc selon les descpitions et critiques d'Ambroise, être passé d'un partage inné des tâches, des devoirs, et du pouvoir, à une vie centré sur l'ambition personnelle et l'accaperement des privilèges, laissant l'idéal politique ambrosien au stade d'utopie inateignable.

\bigskip

La conception du pouvoir politique chez Ambroise, en prenant pour modèle la réussite de son observation du monde animal, ne peut donc pas se résumer à un changement uniquement gouvernemental issu d'une réflexion théorique. Pour y aboutir, Ambroise décrit qu'elle doit venir d'un changement collectif qui se construit autour du rejet du désir de pouvoir, d'une lutte contre la paresse, et plus généralement d'une prise de conscience du travail pour le bien commun. Il cherche à former une responsabilité morale pour revenir à ce qu'il décrit comme l'ordre naturel établi par Dieu, permetant l'établissement d'une véritable république. Ainsi, après avoir clairement décrié les limites de la société romaine, et établi son idéal politique, Ambroise s'attache, dans plusieurs de ses oeuvres, à définir avec plus de précision les vertus et les concepts nécessaire à la bonne gestion du pouvoir et de la société.
